\section{Esoterism and Christian Mysticism}

\begin{quotex}
This essay was originally published by \textbf{EA}, or \textbf{Julius Evola}, under the title ``Esoterismo e Mistica Cristiana" from the third volume of \emph{Introduzione alla Magia}. Obviously, Evola was quite interested in the topic since the essay runs to 20 pages and he was quite familiar with the secondary literature.

Note that Evola's critique is serious and has no similarity with the puerile critiques by some neopagans who pretend to be following his lead. Ultimately, it is unclear what he accomplished, or what he hoped to accomplish. If this critique stands as is, it is a Pyrrhic victory, more like a man sawing off the tree limb he is sitting on than anything constructive. Certainly, there is no ``ancient" tradition that can meet the standards he defines here, and of the non-European traditions, there is currently only one that offers a real possibility.
\end{quotex}

On more than one occasion it was said, in these pages, that Christianity represents a religious system which, if it contains various traditional elements, is nevertheless lacking an esoteric and initiatic counterpart. In comparison with what was typical of ancient, or non-European, traditional principles, this constitutes in a certain way an anomaly and is quite far from corroborating the claim of superiority asserted by the religion that has come to predominate in the West.

As to the traditional elements mentioned, they do not refer so much to Christianity as such, i.e., as the pure evangelical doctrine, but to the corpus of Catholicism, with relation to the symbols, myths, rites, and dogmas by which its orthodoxy is defined. Here more than a few elements also apt to draw from a higher plane and to be catholic in the etymological meaning of the word, that is, universal, assume a merely religious form and validity.

Catholicism lacks an esoterism, because there is no regular elite in its hierarchy, endowed with adequate authority that, in order to be in possession of the corresponding knowledge, is aware of the deepest metaphysical and meta-religious dimension of those elements.\footnote{Related to one of the ``intellectual gifts of the Holy Spirit", to the comprehension of the deep meaning of the symbols and sacred scriptures that Thomas Aquinas mentioned. That remains however a simple announcement, no concrete example of an interpretation of the type being found in Thomas nor in the other authors of the Church. In patristics, e.g., in Origen, one finds the distinction of three meanings of scriptures, corporeal (historical), psychic (moral) and pneumatic (spiritual), that last of which is to be discovered through ``analogy", with the use of a ``spiritual intelligence". Here he proposed to transform the sensible Gospel to the spiritual Gospel", on the basis of the principle that ``the Saviour has willed to make symbols of his own spiritual actions". But in the concrete, for the Old Testament this is reduced to an allegorical interpretation that tends of make it a prefiguration of the New; and in this the last word is always the Christian mystery, where the esoteric interpretation should bring back this mystery—as the particular expression—to a metaphysical, universal, and super-Christian plane. The fact that even these authors touch so many points of Judeo-Christian scriptures containing effectively initiatic elements without even realizing it, proves that it would be difficult for them to recognize the gift of ``gnosis".}

Regarding the experiential side, the rather complex problem of the importance of everything that in Catholicism is rite and sacrament should be confronted, establishing both the similarity and the difference of the religious plane compared to the initiatic plane. Also at the base of a religious tradition, such as the Catholic tradition, there is a spiritual influence and its transmission in an uninterrupted chain through regular and well defined rites. Such continuity in Catholicism is centered in the so-called apostolic succession. In its principle aspect, the transmission is tied to the ordination of priests and bishops. The other sacraments, the first being baptism, are intended to recover, aggregate, and establish the individual in the traditional current, in the middle of the transformation that by virtue of the principle rites would be held to produce in his nature.

Naturally, this structure is characteristic of not only Catholicism, but is found also in every traditional form, even those not specifically initiatic. Only that in Catholicism the claim is more explicit that the rite has the effect of a supernaturalisation and a divinization of human nature\footnote{In this regard, the expressions of the liturgy of Holy Saturday with reference to the baptismal waters are characteristic: ``Deem it worthy, Holy Spirit, to make fruitful with the secret mixture of your divine virtue, this water prepared for the regeneration of men, conceived by sanctification, may emerge from the immaculate womb of this divine font, a celestial race, a renewed creature." [This is my adaptation of the translation from the the 1962 Roman Missal, modified to match Evola's translation, which he apparently took from \textbf{Kremmerz}]}, which otherwise, through the effect of original sin, would be corrupted and vain: whence the idea that the Christian represents spiritual man par excellence in respect to whoever is not such, and whose salvation is only in the mystical Body of Christ, which is the same as saying in the Church as community and chain formed by the rite and carried by the corresponding spiritual influence. Objectively, this claim can only serve for ``internal use" and is lacking every justification. Every regular tradition can propose it with the equal right as Catholicism, because every tradition is likewise formed on the basis of an influence from above, that transforms the naturalistic element of the individual and gives rise to a new current among the forces of the world. In fact, one cannot then see where the Christian shows himself superior to the members of other traditions and has a spiritual dimension that is nonexistent anywhere else.

We need to emphasize, which often is not done, the fact that in Catholic orthodoxy the rite is conceived with the same characteristics of objectivity, independence from sentiment, from ``psychology", and even from the morality that are typical in the magical and initiatic order. That is evident in the rite of baptism, since its effect would be independent of any intention and merit of the baptized (as is quite clear with newborns); so pure in fact that the sacred quality induced in the regularly ordained priest both acquired once and for all and is not even lost in the case or moral unworthiness and also of unbelief. And the same is valid, in terms of principle, for the other rites and sacraments of Catholicism.\footnote{Strictly speaking, in Catholic theory, one even goes beyond the sign, as far as the belief in the magical power of the rite. Regarding that, it is the case, for example, for the power of ``dissolving" in the rite of absolution, that would certainly suspend the karmic law of cause and effect.}

Nevertheless, even when the conditions are present for the real efficacy of the rite, and this does not live through itself on a plane of mere devout fervor and mysticism, even when, therefore, one admits a certain non-human and sacralizing power of the Catholic rite and sacrament, one cannot confuse the order to that which is typical with the initiatic order, and even less can one think that the former can take the place of the latter. Guenon indicated the difference, in this regard, in the following terms: the religious rite propitiates a participation in the supersensible order maintaining however the individual limit, while the initiatic rite would realize that of a super-individual character; the former would aim for the ``salvation" of the soul of the individual, in terms of a prolongation of his individual existence beyond death; the latter would lead instead to true immortality. The most essential difference regards however the presence or the absence of the theistic premise in the concept of the sacred. Everything that is religion, and especially Christian religion, has as limit the idea of a personal God, distinct, as such, from the creature; for it, the idea of a plane in which this distance is abolished by the supreme metaphysical identity is unknown, with which both ``liberation" and initiatic ``awakening" are defined.

In a complete traditional form, religion and initiation are two domains that do not exclude each other but, though remaining firm in their heterogeneity, admit a passage from where the possibility that has religious value to a higher degree, can also assume an initiatic one. But where, as in Catholicism, that does not occur, religious rites appear in a certain way as a useless and misleading parody of initiatic rites, that sometimes almost seems as profanation, while on the experiential side, the highest peak is represented only by mysticism.

Mysticism is one thing, initiation another, this is a point as essential as it is generally unknown. The tendency to reduce the most diverse states to mysticism is widespread. There are certainly cases in which the mystic passes beyond the sphere typical of his path leading to transcendent realizations; but that implies a true metanoia and always represents an exception — even prescinding from the fact that similar realizations indicate, in such conditions, almost always a fragmentary and confused character. Besides, in Christianity mysticism is presented in a characteristic form, constituting itself almost as a closed system, where the noted exceptions are extremely rare.\footnote{After the period of Greek patristics, they are limited almost exclusively to German mystics, starting with Meister Eckhart.} Precisely in Catholicism mysticism is presented as the simple continuation of religious and sacramental experience, and the cognitio Dei experimentalis, the experiential knowledge of God, that would constitute its essence, in spite of the misunderstanding to which one can be led by some expressions, remains in the domain of subjectivity and affectivity, and has little to do with pure intellectuality, with the effective destruction of human nature and real deification.

\paragraph{Esoteric meaning of silence and homelessness}
Anyhow, it will not be useless for our readers to delve deeper into this question, considering at various points the meaning of asceticism and Christian mysticism according to an esoteric valuation.

Above all we have to notice, in the history of Christian mysticism, a process of degradation in a certain way. If we refer here to the first centuries of the Church, we find in patristics not a few elements still having a connection with initiatic teachings, and that due to the simple fact that is was formed in the environment in which the influence of the traces of the preceding and co-existing mystery wisdom and neoplatonism was still alive. From Augustine on, and then with the Spanish mystics, there is a rising humanization or psychologization of Christian mysticism, which not even the Thomist stream avoids. The subject of mystical experience becomes more the soul rather than the spirit, the affective aspect prevails, and moral views take the upper hand over the ontologico-existential.

The fundamental character of Christian mysticism is passivity. According to the Thomist conception ``the fundamental element of mystical contemplation consists in a passive movement of love for God, which often depends on a certain feeling of his presence".\footnote{Padre F. D. Joret, \textit{La contemplation mystique selon s. Tomas d'Aquin.}} It ``is not the product of our activity but rather it is usually preceded by asceticism and meditation, and is normally received as the reward of one's own efforts. But it depends on God's initiative; and our spirit is passive in it, although it reacts vitally under the divine impulse."\footnote{\textit{Ibid}, p. 9} This idea of passivity is accentuated with the theory of double grace, one called ``sufficient", tied to normal Christian life, and the other extraordinary, due to the mystical gifts of the Holy Spirit, who would permit them alone a truly supernatural development. The same arising in the soul of the love of God—which takes the place here of the fire of intellectual initiatic asceticism—is conceived as an effect of divine grace.

Some positive aspects of Christian asceticism, as preparation for contemplation, were present especially at its origins. We will note them.\footnote{On this topic we refer to two books by Father Anselmo Stoltz, \textit{L'ascesi Cristiana} [AC] and \textit{La Teologia della mistica} [TM], part of the current that is seeking to revalue Greek patristics.}

A first guideline of asceticism is isolation and simplification. It, moreover, is justified in terms of analogy, with reference less to Christ than to the Father, and reflects the earlier pre-Christian ideal of ``deification".

``The man, who through the gifts of grace and his continuous esoteric effort will have to arrive at resembling God, must necessarily be an image of the simplicity and unity of God…Little by little deification increases, man approaches always closer to the perfect divine simplicity that asks nothing outside of itself, that is sufficoent in itself. It is obvious that a theology that (like the Greek) gives so important a place to the concept of deification and that conceives God first of all according to the aspect of a single, simple being, sufficient in Himself, will have to see the ideal of Christian perfection in man completely free from all passions and tends to God reproducing in him, as far as he can, unity and simplicity" [AS]

Originally, a not different base had the anchorite as the ideal, physical external isolation, in this case having to have, therefore, only a symbolic and ritual value, hence that of a discipline limited to the training period. In a higher degree, the \textit{extra mundum fieri}, the isolation from the world, is no longer necessary to occur materially, but only in the spirit. In that regard, Christian and initiatic asceticism can meet halfway.\footnote{Cfr. Islam, which has no monastic asceticism, the concept of \textit{zuhd}, which is not actual material renunciation, but inner detachment.}

Ascetic detachment can also assume the form of \textit{instabilitas loci} and \textit{xeniteia}. The detachment from one's own land and from one's family or people, eschewing a fixed home, wandering, bringing oneself to foreign and distant lands—all that can likewise have a symbolic or ritual value in order to bring to awareness the idea that one is not in this world as in a definite place and in one's own country, but rather as a wanderer and an exile, not forgetting one's celestial origin and end.\footnote{AC, 53-60. In such visual persecutions, expulsions, and similar situations, it would be of value as extraordinary means used by God to remind us of what is our true country.} This meaning, as was noted, was very alive in early Buddhism and, we can mention, with a more special connection to the life of action, was not extraneous to the chivalric ideal in the form of the ``knight errant". One must nevertheless note that in Christianity the premise for the coherent realization of a more general meaning is missing, because Christianity rejects the initiatic theory of the pre-existence of the soul to this terrestrial life, and only with reference to this theory can terrestrial existence signify the passage of one who came from far away and again turns towards other states of being.

Christian asceticism also knows the technique of silence. ``That is not only a renunciation of communication with other men, but is also a positive factor of the interior life" [AC, 62]. We know, therefore, the part that the discipline of silence had among the Pythagoreans, as well as the meaning of yogic and initiatic silence in general. In Christian asceticism almost all the practice of the more interiorized degree of such discipline is lacking, that is ``being silent" not only with the spoken word but likewise with thought (Ibn-Arabi's ``not speaking with oneself"). The so-called ``prayer of quiet" is known, where we meet the expression quies contemplationis and Saint Gregory, among the conditions for mystical contemplation, touched on freeing oneself from the noise of thought [TM, 112]. A precise and active method, in this regard, which is known in the East, however is missing in Christian asceticism. And in the monastic forms, as the anchoritic ideal is supplanted by the cenobitic ideal, more consistent with the collectivizing tendency of Christianity, as well as the liturgic element, brought to a hypertrophy, predominates over the individual discipline of silence even in contemplative orders, such as the Carthusians, the Carmelites, and the Camaldolese.

In Christianity in the state of, what we call, ``silence", one is more inclined to follow the techniques than to cultivate certain internal dispositions. The so-called ``cognition of external things" comes back to that.

``It is actually the knowledge of Satan's kingdom, in deriving the power that this exercises on the sensible world and in intuiting the true place that material things must have in relation to the divine plane. This knowledge must not remain purely speculative: it is necessary that the soul at the same time express it in practice, i.e., that it frees itself from the disordered affect for the creature, incompatible with the life of union with God". In such regard there are those who think that such ``knowing" must be received from angels, no natural and human science being able to give it to the soul. [AC, 130-3]

Initiatically, ``separating the real from the unreal" corresponds to that, that are achieved essentially through lives of pure intellectuality, loosened from emotional being; the neutralization of ``demonic" forces is one of its natural consequences. In the second place, always with the same goal, it is said that the soul of the Christian must know what it is before God, which would be tantamount to ``being convinced of his own nothingness", freeing himself from egoism, redeeming himself from his own will". Such a discipline undergoes nevertheless the deformation peculiar to the Christian way of understanding ``mortification" and the ``creaturely" passive attitude. Initiatically, it is about ``disidentification" and surpassing oneself as simple individual being, a surpassing, that does not come about through degradation and humility, nor with reference to the image of God as a distinct being, but rather in conceiving one's own persona as a contingent mask and in itself irrelevant in relation to the true I, almost as in the consciousness of an actor who does not confuse himself with the part that he played. In every way, these are the two principle ways through which, the knowledge of external things of oneself as individual being achieved, Christian asceticism seeks to propitiate the quies contemplationis and to not be more susceptible to inner things, or to passions and frivolous desires.

\paragraph{Sacramental life and initiation}

If we now consider the grades in Christian asceticism beyond the preparatory levels, those which are realized in it and are actually mystical, the fundamental point, which is almost exclusive of references, is the imitation of Christ on a sacramental base. Here we will omit the purely devotional and emotional forms, which lack any value for us; that they were able to take root in active and volitive races, in general, as are the Western races, remains an enigmatic point. The mystical view that is encountered on a slightly higher plane is the following: Christ has transfigured the human nature he assumed by incarnating himself in the body of the Resurrection and Ascension. And the glorification of the body of Christ must be likewise understood in its deepest sense, in relation to the Redemption as restoration and fulfillment of the life of our body in paradise. Christian asceticism considers moreover the event of Christ as the way to glorification with reference to the restoration of the Adamic state in which the body was completely subject to the soul, and its sensory part to the spirit. And since in the redemptive work of Christ the way toward a renewed harmony between the soul and the body is passed through death, and since we must only follow this way in order to assure for our body victory over the effects of original sin, so Christian asceticism has as its central point the concept of mortification, of dying and being reborn in Christ [TM 182,183].

That was accomplished previously in the sacraments beginning with Baptism. It is believed that redemptive and sanctifying grace shared in man by baptism already contains the seed of supernatural mystical life, because one is baptized in the sign of the death of Jesus and the crucifixion of ancient man. If in baptism, the faithful imitate and follows Christ into death, and in the Eucharist the food and nourishment of new life, they would participate mystically at the sacrifice, in the resurrection and ascension: a real participation already in the mystery of the sacrament, even if it is fulfilled in a veiled, invisible way, not accessible to our experience [AC 71-79].

Therefore, the idea prevails in the orthodox Catholic conception, that there is no essential discontinuity between the sacramental life of the common believer and the mystical life: the latter should not be exceptional, it was already contained germinally in the former. The mystic, instead of leaving the seed of union with God inactive, brings its energy into actuality to the point of realizing this union as experience [TM 44-50]. A more concrete interpretation is the following: ritual (sacramental) participation in the death and resurrection of Christ does not transform the entire man in a single blow. It is first of all the interior man (\textit{eso anthropos}, Rom 7:22) who is to be transformed. Then the ascetic, aided by a special grace, will have to kill little by little the ``law of sin" in his limbs and to prepare in that way, according to the new supernatural principle of life received, a new body, the resurrection of the flesh, physical death ultimately to would divest man from the ``body of sin" [TM 39-42]. In terms of dogma, this realization is nevertheless remanded to the time of the ``universal judgment". It is not admitted that one can accomplish that in life.

We now see how one can judge the whole of this doctrine from our point of view. In this area, there is no need to emphasize that the model of death that opens the way to the supernatural order is not specifically Christian; even if it is not essential for every way of spiritual realization (because it is not claimed that this must always entail a crisis), it can have initiatic value. Specific to Catholicism is the idea that it is not an exceptional and dangerous initiatic operation, but the religious rite accessible to everyone that has a power to induce in man a supernatural quality, even if only germinally and potentially, and to produce it in a certain way from the outside, without any participation of the person, even beyond his consciousness, as is evident in the case of infant baptism.

Even if, as was mentioned, one speaks of the further actions of the grace of the Holy Spirit for mystical development, this view contains the principle of a dangerous confusion of planes. If conceived in terms of reality, death, resurrection, and glorification in the sacraments, it cannot be other than symbols and prefigurations, concealing an order that transcends the generic sacralization connected to religious life. And it would be erroneous to suppose that the Christian, having had baptism and participating in the other sacraments of his tradition, finds himself for that reason in any position of advantage through the realizations foretold by an effective initiation. The Alexandrian distinction between the pistikos and pneumatikos, i.e., between the simple believer and the gnostic or initiate, must be maintained. Only through mystical experience can the idea fit, that is does not necessarily imply a rift, but can be an almost natural development of life which, in the Christian, is supposed to be supernaturalized by the work of the sacraments. Sacramental life, as pure mystical life, cannot go beyond a psychological, subjective or moral plane, while the initiatic life has an ontological and super-individual character. In regard to the former one can speak, at most, of sanctity, not of deification.

As for formal correspondences, through some details we can point out that baptism is a simple ritual image of contact with the ``Waters" of ``dissolution" in the life-principle anterior and superior to every individuation or form: experience, that in its radical, exceptional, and even dangerous character cannot obviously have any comparison with whatever is happening —even if invisibly— in Christian baptism.\footnote{The Gospels speak of a baptism by fire beyond that of water. This double baptism and double regeneration corresponds to the two phases of initiatic work called albedo and rubedo in Hermetism.} The notion of the ``glorified body", in which the law of death is conquered, comes directly to Christianity from the earlier mystery traditions, losing nevertheless, in the form of the dogmas of the resurrection of the flesh or of purely eschatological perspectives in the afterlife, its concrete and initiatic meaning. The scheme, actually, is already found in the rite of the \textit{transmutatio} [transmutation] of the Eucharist: the bread counts as the body, the wine as blood and soul. The one and the other principle are transformed and when, after that, a piece of the host is united to the wine, the meaning is the joining of the transfigured soul with the body in the manner proper to the resurrection body or immortal body of Christ.

Now it would be the case of speaking truly of superstitions whenever one presumes, in the Eucharistic participation, anything more than a simple allegory, with effects of moral elevation and, if one prefers, of mystical rapture: certainly not that which always is taken as the furthest limit of possible accomplishment on this earth to an adept. In regard to all this, whoever has a notion of what it actually is about, would come to think of it almost a profanation.\footnote{Cfr. the article on the ``immortal body" in Introduction to Magic.}

Aside from the hoped for intervention of grace, in the purview of realization, i.e., of the development and the actualization of the influences induced by the sacraments, there is no precise way recommended to the Christian. The simple subjective dispositions corresponding to the so-called theological virtues, i.e., faith, hope, and love, remain. ``Mortification" is conceived in essentially moral terms, with undeserved stress on the value of everything that is suffering and ``penance". Here there is actually an incompatibility with what is typical to a healthy and normal Western type of man: on the basis of the sense of ``guilt" and the congenital sinfulness which every man must be conscious of, we know that almost pathological tendencies have often appeared specifically in Christian mysticism. Moreover, prayer, oration, understood as a gift of the Holy Spirit, is conceived as the fundamental power of the mystical life, to the point that ecstasy itself is considered by some as a level of it, the highest level [AC 121].\footnote{The contemporary writer, \textbf{Frithjof Schuon}, who strove to find initiatic prospects in Christianity, believed that the active counterpart of passive participation mediated by the sacraments was (with specific reference to the Eastern Church) the practice of the invocation of the saving name of Christ. This is one of the more primitive techniques for the ``killing of the manas", i.e., for the neutralization of the mental ego, analogous to the continuous repetition of the ``divine Name" in Islam. Moreover, the abuse of the liturgy in the Catholic contemplative orders does not have objectively any other end. That, beyond the negative aspect, the ``virtue of the name" also enables certain states of illumination, is something problematic. In every way, it is a chance event, and it is necessary to have faith in the protective actions of the sacraments in order to unwrap the possibility of an action of the most diverse extra-sensible influences, once the manas is ``killed". In any case, the mystic has little way of interpreting correctly the phenomena that occur, because it is held on an emotive, rather than intellectual, level, and the Christian devotional framework, with its various images, rather than helping him, serves instead to lead him astray.}

\paragraph{The gifts of the Holy Spirit, the tollhouses after death, early Christianity and mystery cults.}

According to \textbf{Pseudo-Dionysius the Areopagite}, the fundamental phases of Christian mysticism were the purgative, illuminative, and the unitive ways, a schema which, in that form, can be also be found in the initiatic path. As a reality, however, things were quite different. Therefore, concerning the first phase, Guenon's comment is absolutely correct, when he points out the confusion due to the use of the terms ``union" and ``unitive way" by mystics.

\begin{quotex}
This union does not have the same meaning as Yoga or its equivalents, so that there is only a wholly outward similarity. It is not illegitimate to use the same word, because even in current language, one speaks of union between beings in many different cases where there is obviously no degree of identification between them, but it is necessary to take the greatest care not to confuse different things under the pretense that a single word is used to designate them both.\footnote{Guenon, Rene: Direct contemplation and reflected contemplation, from \textit{Initiation and Spiritual Realization}.} 

\end{quotex}
Moreover, the frequency of the use of the symbolism of matrimony for the mystical union is significant. It does not mean a unity that maintains the distinction that occurs in uniting a man and a woman, in the same way that the relationship of an ``I" and a ``Thou" exists in prayer, even in its highest levels. Only by analogy can one speak of a cognitio Dei experimentalis—experiential knowledge of God—in Christian mysticism. A condition of ``eccentricity" (in the sense of non-centrality) is essential, almost without exception. That, moreover, is indicated by the term ``ecstasy" which means ``to go out from oneself". Initiatically, it is more about the opposite, i.e., of a coming back into oneself, of making oneself ``central", through which someone correctly spoke of enstasy, instead of ``ecstasy" for the states of Yoga.

The phase of the illuminative way could induced by putting it in relation with the effects of the gifts of the intellect of the Holy Spirit that are six [sic] in the Thomist conception: namely, to be able to discover

\begin{enumerate}
\item the substance behind accidents; 
\item the inner meaning behind the words of sacred texts; 
\item the truth behind the symbol; 
\item the spiritual reality behind the sensible types, i.e., behind the phenomena which appear to common experience, and consequently also the hidden effects of a cause and the hidden cause of specific effects. [Joret, Op. cit., p 274] 
\end{enumerate}
Formally, it would be a question of powers of an effectively initiatic character, connected more or less to ``intellectual intuition". But it suffices to read the same examples made by St Thomas in order to see that there is none of that, or rather that is does not go beyond a theoretic framework.

Concerning the purgative way, again with the goal of indicating the reflections of initiatic views, let us reproduce this citation:

\begin{quotex}
Many writers of Christian antiquity affirm that the souls, when leaving this world, must pass through various tollhouses of the demons of the aerial region. They examine the soul to see if they can find something in it that belongs to them … The soul is not able to continue its voyage until it has been purified from all its dross. At the same time, beyond purification, the soul is also instructed by good angels in every divine science, so that it is capable of understanding the Mysteries of God. This whole process of purification and instruction is repeated until the soul has arrived at perfect purity and the fullness of knowledge. This gradual purification after death generally is no longer admitted by Catholic theologians in all its particulars. It was however practical to the anticipation of our asceticism that is confirmed in the life of prayer, and in such applications, the aforementioned doctrine expresses some absolutely indisputable truth. [AC, 128-129] 

\end{quotex}
Now, in all that, the reflection of the mystery doctrine is quite visible, such as the Mithraic, with the passage through the septenary or planetary hierarchy, which does not have to do with the world of prayer, but with that of initiation, in which, as the same ``purifications", were often distinguished by exactly seven grades.

Other ``residues" not lacking interest are the views that, in the ancient Christian writings, were about the restoration of the Adamic state, understood as prior to the goal of Christian regeneration, because in patristics they insist exactly on the continuity on the fact that one must achieve the perfection characteristic of the primordial state, by embarking on the search for the lost terrestrial paradise.

It was the ancient Christian idea, maintained moreover up until the 16th century, that this place, from which our first parents were driven away, still exists, in a higher region, inaccessible to men, unless they are helped by exceptional divine grace, as would have been the case for Elijah, Enoch, and even Saint Paul (II Cor 12:1-5). This inaccessibility was expressed with the symbolism of desolate and impassable places, equivalent to that of the angel on watch with his flaming sword or a fire zone that encircles paradise. That copies the initiatic doctrine of the ``center" in its relation with the primordial state: always present from the metaphysical point of view. And the symbolism of passing through the wall of fire is not different from that of baptism of fire, referring in that way also to the theme of Hercules who achieved immortality (he ascends Olympus to marry Hebe, the perennial youth) only after the fire consumed his human nature on Mount Oeta: a theme that refers to still many others of initiatic origin, e.g., that of a virgin enclosed in a circle of fire, who will belong to whoever can pass through it without perishing.

Obviously, in projecting it to an afterlife, the Catholic conception of ``purgatory" reflects the most general theme of purification.\footnote{It is said: Only who is purified through fire can enter paradise. St. Ambrose (In Ps. 118, serm. XX,12)} Moreover, there is also an allusion to the initiatic idea that paradise and the kingdom of Heaven represent two distinct ``places".\footnote{TM, 25-6. Where this passage from St. Ambrose is cited: `Paradise is a permanent region of heaven, but it is, so to speak, the ground floor of the kingdom, the foundation on which the kingdom of heavens properly called is constructed at the summit; it is the lower region of the invisible heaven, from which the elect, each according to his own merits, will ascend sooner or later toward the different higher regions or heavens." In Christian terms, the distinction is also expressed by saying that the grace of Christ is superior to the grace of Adam and leads to a perfection that is beyond the ordinary paradisiac state.}

In the ancient Christian view, there was always a mingling with magic, when, once admitted that the Adamic state can be re-achieved through the other levels of mystical contemplation, they indicate the prerogatives of that state. The first is the gift of knowledge, or gnosis—through which the fact is recalled that Adam had given the animals ``their names"—the name, according to the ancient conception, expressing the very existence of a thing or a being. In relation to that, St. Thomas also writes:

\begin{quotex}
This rectitude of man in his original and divine institution consisted in the fact that the lower part of his being was subjugated to the higher part and this was not obstructed by the other. Therefore the first man did not find in exterior things an obstacle to pure and steady contemplation of intelligible effects, perceived by means of an irradiation of the first truth. [TM, 89] 

\end{quotex}
That, however, is also related to the second gift, which is apatheia in the most general sense of being undisturbed in respect to passion, of not being preoccupied by impulses, and of a natural sovereignty over instincts.\footnote{St Augustine's (De civ. Dei, XIV, 15) view is of interest, according to which at the very point in which the soul separated from the divine world, the body ceased in its turn to be subject to the soul and to obey it, and came under the state of passivity defined by concupiscence.}


\paragraph{Cult of Mary and the initiation of the Fedeli d'Amore}

Moreover, according to patristics, man as microcosm comprises the entire creation, through which the dominating power possessed by man over himself is manifested likewise in the macrocosm: the dominion of man over nature is the prize of the complete victory that man brings back on himself: the hierarchical relationship established between the powers interior to man are transferred into the relationship between man and the world [TM, 84]. It is through such a way that an unlimited power over creation had been promised to the saints, according to the words of Mark 16:17-18.

\begin{quotex}
And these signs shall follow them that believe: In my name they shall cast out devils: they shall speak with new tongues. They shall take up serpents; and if they shall drink any deadly thing, it shall not hurt them: they shall lay their hands upon the sick, and they shall recover. 

\end{quotex}
It is unnecessary to point out that the premises just noted are the same as in high magic or theurgy: Agrippa and Paracelsus did not express themselves any differently. Finally, as the third prerogative of the adamic state, there is the gift of immortality, athanasia, ``corporeal immortality being a sign and proof of the interior presence of God". For that reason, ``a restoration of the primordial state would also include the reacquisition of immortality". [TM 82-4]

At the origins of Christianity, there was, as an echo, the ontological ideal of deification—assimilation to God and participation in his nature—as an effect of the earlier mystery tradition and it still predominated over the morality of faith, love, and the merits that carry essentially their fruits into the afterlife. Therefore, there were still half-closed perspectives of the type that, once the dross has been separated, tends toward a world higher than that of mysticism, and of which the world of mysticism is only a weakened and rather humanized echo.

One final point. We know the part that Mariolatry, i.e., the cult of the Virgin as ``Mother of God", has in Christianity. If this cult, in its most external aspects and from the historic point of view, gives evidence of the influence that gynocratic views had in Christianity, i.e., the pre-eminence of the feminine divinized principle (Magna Mater, Gaia, etc.), as in the archaic, pre-Indo-European Mediterranean cycle, then only an esoteric assumption of it is possible. A rather unique thing, it is suggested by the same Catholic author, to whom we referred repeatedly, when he says: Mary is ``the highest ideal of the ascetic", because ``the ascetic tends to form Christ in the soul, to make it become the Christ bearer and the Mother of Christ" [AC 192].

Although this symbolism, itself attributable to the scope of initiation, at its base contradicts the principle of passivity characteristic of the Christian mystical way and as emphasized by the conception of the indispensability of grace (i.e., of the activity of a force that psychologically is experienced in terms of grace) for every supernaturally efficacious effort. If in Christian symbolism the Virgin is only fecundated—by the Holy Spirit—according to the initiatic symbolism the ``virgin birth", it is, more precisely, the birth that has no need of external help, the endogenesis due to a pure and intact force. Kumari, the Virgin, is the ``power of the will", according to a Hindu text.\footnote{We note the correspondence with the expression used by Dante (Par. XXXIII, 34) for the Virgin: ``Queen, who can do whatever you will".}

Another mystical interpretation of the ``Mother of Christ" is that which makes her the representation of the Church, the mother of supernatural life. [AC, 193] That is, in relation to the supernaturalizing influence that would be connected to its tradition, to its organization, and its rites.\footnote{As much to the form, or doctrine, as to the organization that is its deposit, the symbolism of Woman was applied in various cases in the initiatic domain. So, in the case of the Medieval chivalric initiation of the ``Fedeli d'Amore", Woman represented both ``Sacred Wisdom" and one or the other organizations of that current. The passage at the active movement was often portrayed through the symbolism of marriage, or even of incest (the Woman or the Mother becomes the spouse of the initiate).} Usually, such a point is, however, in connected with the exclusivist claim, a characteristic of Catholicism.

Mysticism considered as a development of the experience of the insertion of the individual, effectuated by the sacraments and above all by the Eucharist, into the current of divine life, and the Church, understood as the mystical body of Christ, as the bearer of this current, is the obvious conclusion that there cannot be true mystical life if not in the Church [TM 77]; and that Non-Christian mysticism, being less than union with Christ (which, evidently, is considered along the lines of an historic personality conceived as the Redeemer, not of a superhistoric and universal role model to imitate), is in fact naturalistic, if not the result of demonic influences. In such a way, we end up presuming that in the non-Christian Mysteries, there could not have been any question of true deification, but only of more or less contingent psychic states. [61-71]

The well-known principle Extra ecclesiam nulla sulus [no salvation outside the church] sanctions, in general, the idea that only Christian redemption could result in liberation from diabolic power, expressed as the breaking of the circle of the demonic force that surround us, while helping us to doubt its power, and that making us think that redemption in the sign of Christ and the Church is superfluous, would be the most dangerous victory that Satan can attribute of his activity [66].

We previously said that all that can only have value for Catholic ``internal use"; likewise, we noted that these views have an exclusively practical justification, no different from that of exclusively analogous views that are encountered in outer, exoteric forms of most other traditions: but they still lack every foundation beyond the sphere of the jurisdiction of each tradition.

In relation to that, we can pose, as the last issue, the question of the extent in which, for those who have an initiatic vocation, it is useful to adhere to and join a religious tradition—specifically, to Catholicism in the case of a Westerner.

As the conclusion of the preceding examination, it seems clear the Christianity, as compared to other traditions, has a particular character. In fact, from one side it is not, like ancient Judaism or orthodox Islam, a pure religion of the ``Law", but that inner experience is important to it. On the other hand, it does not know the experiential, esoteric, and initiatic plane, so that its level is inferior to that of traditions where this plane is adequately taken into consideration. This intermediate nature of Christianity can be characterized by defining it exactly as an essentially mystical religion, which has absorbed and adapted some esoteric elements into its mysterico-sacramental form.

Now, a current that is defined in such terms has a predominantly psychic and collective, rather than a spiritual and metaphysical, character. In every way, that seems rather noticeably what Catholic tradition has become in our time, so that it leads us to think of one of the cases of entities, from which, through involution, the influences that are truly from above have in a good measure receded. But since psychic entities of that type conserved their vitality and strength through inertia, uniting itself to it, for those who have an initiatic vocation, it can serve more of a bond than as the basis for higher development. We come back to a ``chain", from which it is difficult to free oneself and whose subtle influences are difficult to control. So through the doctrinal side, there is something characteristic that even those who sincerely live out the Christian faith sacramentally and mystically are the most ``taken", the most fanatic, the most incapable of recognizing and respecting anything which is anterior and superior to their tradition and that has been realized in other equally legitimate forms in the world and in history.

In every respect, in the case of Christianity, it must be judged that the relations between the outer religious way and the metaphysical way are minimal and that one can immerse oneself in the mystical-devotional current to the point of reaching relatively high levels without perceiving anything of the initiatic and metaphysical order. With its character of a closed system in its own way, with elements which represent only a formal, reflected image of the mystery of transformation and of deification, we would almost say ``lunar", Christianity is perhaps, among the traditional forms, the one which is least recommended to those who want to enter the ``direct way".

\flrightit{Posted on 2013-05-19 by Aeneas }

\begin{center}* * *\end{center}

\begin{footnotesize}\begin{sffamily}

\texttt{X on 2013-05-19 at 18:53 said: }

One of the most valuable articles on this site.

Such Evola is sane, not that one who talks about Nordics and Jews all the time.


\hfill

\texttt{Ash on 2013-05-19 at 20:09 said: }

As I understand it, Evola is saying that the Traditional teachings can use the Christian and Catholic faith as a language to explain themselves, but that this language has shortcomings (elements of the sacraments, the personal God, etc). Would that be correct? It also makes sense of what he says about those revolting against the modern world: ``When we appear to be destroying, we are in fact rearranging and replacing what is on the wane with higher forms, forms more vibrant and glorious." Might one say, he is seeking a new and better language?

If so, that would be a rather heartening clarification of my understanding of the interrelation between Tradition proper and the faiths in which it manifests itself. The language of the Triune God has of course been used to great benefit in leading people to truth, but to his literal existence I have usually taken the point of view of Buddhism or the Vedanta. I rather hope we'll see Evola touch upon the person of Christ. While the Logos is central to the students of Tradition, the teacher himself as a historical figure has been less discussed in relation to the more esoteric currents within his Church.


\hfill

\texttt{Mihai on 2013-05-20 at 03:38 said: }

There is much to combat in this article (which I am sure I have one or two years ago integrally). Right now, I do not have the time, but I will limit myself to saying this:

1. The so-called barrier between esoterism and exoterism is never as absolute and strict as it is usually claimed today. The opposite is usually true, since the degree of esoterism and exoterism are usually determined by the capacity of understanding of the one receiving these doctrines. Nor should such a schema by regarded as universally true, since it is not true in Christianity, where esoterism is the upward extension of dogma. (Of course, I speak from the point of view of human realization- the dogmas are actually the downward extensions of esoteric principles in reality).

2. There is also this procroustean bed of delimiting esoterism and mysticism. This, of course, has its origin in Guenon. However, let it be said that in the Christian East, mysticism has nothing to do with anything of a sentimental and individualistic nature as it is regarded in the West, and, from this point of view, mysticism here is the same as what Evola and Guenon mean by `esoterism'. There is however, a problem in these authors' conception of what constitutes actual esoterism and what ``mere mysticism" or ``mere theology". For example, I don't see why a St Dionysius the Areopagite, Maxim the Confessor or Simeon the New Theologian should be all regarded as mere ``exoteric theology", while the writings of an al-Ghazzali or ibn-Arabi as pure esoteric teaching. What exactly is the difference between these authors, even when regarding expression ?

It is clear that the fault here is to be found in the too narrow conception of esoterism as having to do with secret rites, completely separated from the exoteric dogma. When discarding this all too procroustean ideas, the truth about Christianity will appear in its true light. 

There is one phrase here that is completely unwarranted :

``religious rites appear in a certain way as a useless and misleading parody of initiatic rites, that sometimes almost seems as profanation, "

Even if I were to accept as true Evola's views on Christian esoterism (or the lack of it), I don't see how an incomplete understanding can be a ``parody or a profanation" . To understand incompletely means to not derive the full efficacy or benefices from a doctrine. A parody is the subversion of a doctrine for nefarious uses.

Evola is so focused on what he calls esoterism and initiation that he goes so far as to think that everything religious is simply useless or of limited appeal and only good enough for the masses. 

My question is: when one strives for the highest peaks, doesn't that same one have to, first, climb through the simple roads at the base, before reaching the barren cliffs at the top ?

To strive for a higher realization does not mean to reject what the masses are doing, but to surpass it. If one cannot even do as much as the ``mere exoteric" does, how can one legitimately have pretenses to esoterism ?


\hfill

\texttt{Graham on 2013-05-20 at 20:57 said: }

``the effective destruction of human nature and real deification"

This stood out to me. ``Real deification." The Christian theologians I've read, the few who mentioned the possibility of deification or divinization, always take care to clarify and limit, to stipulate in what sense deification is achievable, in a way that the traditionalists don't. I might summarize the difference as `participation' versus `identification'. It could be a difference in perspective, theological speculation and metaphysical realization, or (for argument's sake) a limitation in Christian theism itself, as Evola maintains. I don't know. I do see it as a stumbling block in the Gornahoor project, when the Christian deification is equated with Sankara's ``atman is brahman"; it seems facile, but I could be wrong. The recent entry on the Thomistic doctrine of Personality was stimulating.

Now to the other part of that sentence that struck me, ``effective destruction of human nature." This is characteristic of Evola's style, and it's what the younger just-post-Nietzschean me found attractive about him, this pungent contempt bordering on hatred for the human. Today it strikes me, distinctly, as heretical, perhaps Nestorian or Apollinarian, and definitely smacking of false Gnosticism. Christ himself has human nature. That Evola would choose to speak of the ``destruction" of human nature, instead of its purification through grace, to my mind suggests a malady. Correct me if I'm wrong.


\hfill

\texttt{Graham on 2013-05-20 at 21:15 said: }

To clarify, I was speaking of heresy in a `pan-traditional' sense, and attempting to apply the specific heresies mentioned in a metaphysical sense with respect to deification, rather than with respect to the historical person of Christ. So if there's a heresy that holds that Christ had no human nature, but only a divine nature coupled with a human body, this Christology has obvious implications for deification; on the other hand, the Catholic doctrine holding that Christ had both a divine and a perfect human nature suggests a deification which Evola deems somehow insufficient.


\hfill

\texttt{August on 2013-05-21 at 00:32 said: }

One might object that there can be no heresies in a metaphysical sense, within which domain one is free to exert his will to the extent of his capability, up to `identification' if one's power allows it. This seems to be the spiritual attitude of the warrior that is often ascribed to Evola.

Specifically, this attitude might smack of heresy because it suggests that God, or Pure Being if you like, is not something entirely fixed and active, and is only a certain `domain'. within which manifold beings aspire to take possession of a great power. This is problematic for those who conceive of Being as culminating in a personal, theistic way.


\hfill


\texttt{Mihai on 2013-05-21 at 03:31 said: }

``is not something entirely fixed and active, and is only a certain `domain'. within which manifold beings aspire to take possession of a great power. This is problematic for those who conceive of Being as culminating in a personal, theistic way. "

This is plainly luciferian non-sense, regardless if we conceive the Absolute as personal or impersonal. Read the Hindu Upanisads- in no way do they suggest that Brahma can be conquered through individualistic will-power and a Promethean attitude.

Come to think about it, not even the gnostic heresies or any kind of heresies for that matter claim this. The only possible exception to this may be some of the most degenerate tantric sects. 

On the other hand, I really have come to ask this, since all too many traditionalists take this for granted : why is an impersonal conception of Deity ``clearly" superior to the personal one ?

``This is characteristic of Evola's style, and it's what the younger just-post-Nietzschean me found attractive about him"

Yes, same here. And it is precisely this reason why Evola is such an excellent bridge between the abyss of nihilism and the plains of Tradition.


\hfill

\texttt{August on 2013-05-21 at 06:03 said: }

You seem to misunderstand slightly, yet your critical attitude is welcome, so I will try to explain better.

There is no question of conquering Gods, let alone doing so through individual will. Rather, it is about achieving an equivalent existential status, and thereby establishing oneself (one's Personality) as the center of a whole state of being. What else is meant by complete freedom in the context of initiation? For this to be possible, the principle of Being cannot be equal to Brahma, or the particular God of any tradition, so long as that tradition ascribes any characteristics to God other than pure Being, which is the potentiality of the indefinite possibilities of formal states. Hence, beyond Brahma there is the Supreme Unity, Brahman, because Being is not the Infinite.

As for me, I do not necessarily consider an impersonal conception of the Deity as superior to the personal; the reason I consider them both is because both must have their place. If the impersonal conception were senseless, or impossible, there would be no need to consider it as it would not be true. If it is true, it must be taken into account, and an objection based on preference does not free us from this fact.


\hfill

\texttt{Ash on 2013-05-22 at 15:52 said: }

Perhaps we should consider the effect having a personal theistic God would have on the formulation of a doctrine of apotheosis. In my mind, it would probably be heretical to say that the human ego is identified with the persons of the trinity or their essence. It would not be to say that the divine element of Being that is the foundation of all states of existence, including human existence, is identifiable with the Godhead. That is, not the person of God but rather the transcendent reality which the personal God puts a face on. A christian would need to differentiate between these meanings, leading to different definitions of words like ``identification."


\hfill

\texttt{Cologero on 2018-05-19 at 21:32 said: }

There is a way out of the impasse. The first step is to recognize a knowledge path within Catholicism. We are commanded to ``know, love, and serve God", i.e., the three traditional paths of jnana, bhakti, and karma yoga. Pace Evola, there is a path of divine union, not just the duality he mentions. E.g., St. Catherine of Genoa hints that the Self is God.

Moreover, Gornahoor has followed the suggestion in note [4]. That accounts for the emphasis on the Fathers (via the Philokalia) and the stream inaugurated by Eckhart, in particular, Ruysbroeck and Jacob Boehme.


\hfill

\texttt{Jason-Adam on 2013-05-26 at 21:16 said: }

Catholicism is incompatible with what is ``healthy and normal" for the ``active and volitive" Western man……and what is the source for Evola's characterising Western man this way ? Nietzsche doesn't count as the man had no transcendent knowledge so why repeat his prejudices ?

The dismissal of Schuon strikes me as very superficial. I've spoken to hesychasts and they told me that actively repeating the Jesus prayer does induce a transcendent state, it's not a mantra. And also the hesychasts have methods for recognising false and deceitful visions.


\hfill

\texttt{Cologero on 2013-05-26 at 21:56 said: }

Well, Jason-Adam, as Guenon pointed out in a letter to Guido de Giorgio, Evola was full of prejudices. That is unfortunate, given his otherwise brilliant insights. Evola, apparently, was not so concerned with intellectual coherence. For example, he referred to the time before the French Revolution as ``healthy and normal". Were the medieval knights not healthy and normal? Painting with such broad strokes does not produce any art.

Even if the hesychasts are at a ``low level", that is still much higher than anyone else. Evola included, since he was never initiated into anything.


\hfill

\texttt{Mihai on 2013-05-27 at 03:21 said: }

The problem, it seems to me, is the terminology employed. Since the Eastern Church (as did the One Church in the beginning of Christianity) describes its methods, disciplines and knowledge as mystical, and Evola identifies ``mysticism" solely with what came to prevail in the West under this name- that is individualist and spontaneous achievements of an admittedly sentimental and subjective nature- he comes to the conclusion that the Jesus prayer must be nothing more above this sentimental and subjective experience.

He clearly had no direct insight into the mystical teachings that have been kept here until the modern times, he probably had no adequate knowledge of the Philokalia- and this is proven by the fact that he doesn't once cite a given Church Father, but only modern authors talking about the Fathers, so his conclusions in this area are terribly superficial. 

Mysticism means here something that resembles his definition of esoterism- and I say resembles, because even his definition of esoterism is somewhat strained and certainly only includes a portion, not the whole of what used to be `named esoteric' knowledge throughout history.


\hfill

\texttt{Mihai on 2013-05-27 at 03:27 said: }

``Even if the hesychasts are at a ``low level", that is still much higher than anyone else. "

True, and I would even go further to say that even doing nothing but the `external' discipline of going to church and receiving the Sacraments is still better than doing nothing but stand in front of the computer screen and complaining on `traditional' forums of how the West, because of Christianity, has lost all contact with esoteric knowledge- much like the pseudo-traditionalists of our days are doing (not surprisingly, many being affiliated to the new-right).


\hfill

\texttt{Logres on 2013-05-27 at 13:51 said: }

``The Alexandrian distinction between the pistikos and pneumatikos, i.e., between the simple believer and the gnostic or initiate, must be maintained."

Certainly agree with this; lately, reading Ousepensky's The Fourth Way, and certainly (for myself, and perhaps other believers) would agree that there is a danger in Christianity today of assuming that one ``knows" when nothing of the kind has happened, when the individual has merely transferred illusions to another plane. If there is a danger of giving in to a ``higher" illusion, that danger can be obviated by efforts not to succumb to further illusions. The danger, inherent if any, could come merely from institutionalizing those tendencies to illusion in the forms and sacraments, which certainly goes on, but can still be obviated by the individual, if they are being guided in that Tradition by a higher spiritual impulse. There are many resources, within the Faith, to aid in that possibility. I don't think blaming Catholicism or Christianity in toto, as Evola dances about the possibility of doing, flirting with it, so to speak, serves a constructive purpose in and of itself.


\hfill

\texttt{Logres on 2013-05-27 at 13:54 said: }

F. Seraphim of Sarov, for example, is spoken of in the Eastern Church, as actually having transcended monasticism, and gone back into the world. Isn't there more than a ``trace", here, of recognizing that the journey doesn't end in mystical absorption into God?


\hfill

\texttt{Jason-Adam on 2013-05-27 at 17:01 said: }

Evola strikes me sometimes as getting close to schizophrenic – I mean here's this guy who's an admirer of De Maistre, Cortes, De Poncins, who says the pre-1789 world is the healthy one, yet then he's also into Dadaism, Beatniks, Crowley, is anyone seeing a split mind ?


\hfill

\texttt{Mihai on 2013-05-28 at 03:23 said: }

Not really, but, admittedly, he does have problems identifying genuine spiritual knowledge from counterfeits. For example, he cherry-picks some interesting bits from Crowley found here and there in his doctrine, but he doesn't see the big picture in all of this: namely that Crowley was, in the end, just another occultist whose system is no more than pseudo-spiritual syncretism with certainly lots of material that strikes me as of sinister origin.

(By the way, when I enter the facebook account for Gornahoor, I see a Crowley picture at the top: I believe it harms the mission that Gornahoor has set for itself and attracts all sorts of unwanted individuals). 

In his first volume of the Introduction to magic, there is some material written by antroposophists, which (the material) is certainly filled with occultist concepts.

I think Evola was one of those individuals who feared getting past and leaving behind certain initial steps (such as Nietzsche) of fear of contradicting himself.


\hfill

\texttt{scardanelli on 2013-05-28 at 12:53 said: }

Well said. Your response, I think, sums up the entirety of Evola's problematic relationship with Christianity. It seems that both he and Guenon largely ignored the Orthdox church, and i've often wondered why this is. 

There seems to be competing ideas of Orthodoxy as Eastern Orthodoxy, meaning a cultural phenomenon, and Orthodoxy, meaning the true, historical, ancient form of Christianity. Perhaps the former conception made it possible for them to overlook Orthodoxy. The more I learn, however, the more I tend towards the latter, though I do feel a certain sense of loyalty to Catholocism.


\hfill

\texttt{August on 2013-05-28 at 16:49 said: }

Evola considered the world to have suffered a schizoid event, namely the consolidation of modernity as negation of existing traditions without inauguration of a new one. He did not consider this to be a mere cultural change, but a `concrete' event of cyclical transition that affected the world and man on many planes; not absolutely, but deeply enough that only `esoteric' penetration offered a real escape, whilst even religions had been more or less tainted. Much of this, of course, derived from Guenon, and summarized by Nietzsche as the `death of God'.

Evola thought that the gravity of the change was such that new, modern man could no longer be accommodated by the structures of yesteryear, hence his intrepidity in exploring late modern expressions of this novel atmosphere and its resultant existential turmoil, whilst recognising the fundamental health and normality of the prior traditional forms and the absolute principles they protected.


\hfill

\texttt{Jason-Adam on 2013-05-28 at 17:38 said: }

August your explanation of Evola's thought is very good and I agree with your analysis. 

Do you think it is possible though for someone to make himself a man of a prior age ?


\hfill

\texttt{August on 2013-05-29 at 03:13 said: }

Hi Jason-Adam, here is my comment, feel free to explain your question further if my interpretation is off.

A movement toward his primordial center may be undertaken by a man, along which path he may certainly recover the memory of what it was to belong to a prior age, particularly the essential part, and those ages that are to him directly ancestral.

Yet, this is not the same as making himself that prior man, for his point of departure is inevitably the modern form, that very manhood that could not but emerge after the cosmic shift described above; at least some part of this modernity, even if only the memory, will likely remain a part of his soul.

In any case, the recovery of primordiality is an internal movement, and while its effects may be reflected on the outer man, the individual as such remains what he is within the world in which he manifested.


\hfill

\texttt{Jason-Adam on 2013-05-29 at 14:16 said: }

I can only speak for myself but within me I have always felt that I do belong in this world and I want to live as far apart from it as I can; the more I come into contact with the degeneracy that is modern life I get very angry and want to destroy it all. That is why, before I became aware of Tradition, I was attracted to punk – the pure hardcore nihilist punk to be exact – but since becoming aware of Mediaeval civilisation I do feel an instinctive attraction to it, I feel as if I belong in the Middle Ages. When I first read old works of chivalry, I immediately desired to become a knight, a desire that has not left me at all. That's why I feel alienated from most Traditionalists equally as I do from moderns because they do not share my values and beliefs, which are those of the crusaders. I do feel at home around SSPX circles but those people tend to be of a more intellectual bent whereas I want to take those beliefs and fight for them. Can I be the last Christian warrior ?


\hfill

\texttt{Jason-Adam on 2013-05-31 at 13:54 said: }

Evola seems to be engaging in nit picking over words – what he should be doing is actually trying out both yoga and Christian mysticism to experience for himself the differences and similarities.

The part about the gifts of the Holy Spirit according to St Thomas is the clincher of this piece – Evola admits that theoretically St Thomas is speaking about initiation. We can ignore the judgements he makes afterwords but aside from a personal prejudice I see no reason why one can not use Aquinas as a pathway to realisation. 

Another issue is Evola is so opposed to any notion of a God apart from the Self, but that is a belief not only to Christians but also to Muslims. Evola's self-focus seems to me to be based on modern teachings. It sounds similar to Crowley.


\hfill

\texttt{Mihai on 2013-06-01 at 15:50 said: }

``for the simple (though disturbing to religious exclusivists) reason that a part of him quite clearly saw beyond the specific theology and form of God in Christianity."

Actually, in regards to Christianity, Evola never saw beyond anything, because, like I mentioned in the comments to a previous article, he was totally unaware of the most profound doctrines of Christian mysticism, as preserved in the East. At best, he considered some of their points here and there, but he considered them as somewhat ``copied" from other more ancient mystery traditions, as if Christianity were no more than some fragmentary bits collected from here and there and artificially put together.

In other words, his view of Christianity was quite profane, and the fact that he only cites for references some book of a modern scholar, and never some actual material of Patristic writings, proves that he merely wanted to be done with the thing as quickly as possible, since it never constituted any interest to him.

Regarding some other parts of your post, it is really a good time to put into question an issue which many people take for granted (serious seekers, not only pseudo-esoterists): namely that an impersonal Absolute is somehow superior to the Personal. This is really a mere arbitrary postulate without any serious support, since ``impersonal' is a negation of the person, making it something less, not something more. 

Anyway, both these conception of the Absolute, and Evola's understanding of ``genuine" initiation is an attempt to fit all spirituality into the procroustean bed of what represents the view of just some spiritual ``schools", mainly the Vedanta.

The same point can be made, in some measure, about Guenon as well.


\hfill

\texttt{Mihai on 2013-06-02 at 16:20 said: }

There is here something clearly contrary to any genuine esoterism, and contradictory to what they have, otherwise, tried to point out: namely that Truth is ineffable.

Guenon wrote somewhere that logical relationships are inseperable from ontological ones. Which means that abstract logic can in some measure point us the way as long as we are on the plane of Being. Beyond-Being points to something which is beyond all ontology and, as such, and with which no logical methodology is co-extensive. You can guess from this how coherent Guenon is on this score. 

As such the rationalistic schema above is completely null and void because it tries to fit the Infinite into the limited conceptions of human understanding. It seems to me that, to the traditionalists, Truth must adhere to some verbally constructed schema. Also, I sometimes really ask myself: since it is so easy to talk about the Absolute in logical terms and since the human mind alone can reach and understand in such simplistic concepts the ultimate realities, I really am not sure what we have symbols, myths and mystical methods for. 

This, of course, has nothing to do with genuine apophatism, nor esoterism or higher wisdom, but is only the result of modern prejudice which always feels the need for a radical simplification, in order that the individual may feel comfortable and secure, with a formulation that the mind by itself is able to grasp.

In the Christian East apophatism is taken to a whole new, dynamic level. In complete adherence to the ineffability of God, the Patristic writings do not attempt to resolve the supreme Mystery of God on the plane of human reason and logic. In keeping true to the apophatic spirit and to genuine esoterism, or mysticism, the only claims they make of God are antinomian, mutually exclusive formulations: God is Three and One, the Three are equal, yet the Father is supreme etc. It is here a Mystery which transcends all notions of human coherence and before which the human mind finds itself on slippery ground, completely impotent before the Ineffable- and this is the kind of experience which human weakness detests, hence the tendency of simplistic rationalizing.

If we look at the history of different heresies regarding the Trinity, we can see at their origin the same tendency- the attempt to reduce the real Mystery to the simplistic conceptions of human logic.

The same goes to the Christian idea concerning Divine Personality, which should be understood not as a limitation, but its very opposite.

(By the way, Guenon's claims that the Absolute manifests because it is a possibility, so it simply must manifest. I really wonder what ``kind" of Infinite we are talking about here, when manifestation simply bursts forth from it mechanically, without it having any control and certainly without any intelligence involved. Such a notion of the Absolute is even inferior to what we humans are capable of). 

Besides this, I see in this insistence on impersonalism and the optionality of ``adherence" to the personal, Self-revealed God, which is somehow inferior compared to some verbally constructed abstraction, not the slightest trace of any perennial wisdom but actually another manifestation of modern prejudice, where human personality is dissolved into the atomism of an impersonal and mechanistic universe and where the individual has the satisfaction of holding on to this individuality all the while claiming that he is above all paths and doctrines, because he sees ``beyond" such ``petty human concerns" and therefore succeeds in getting rid of the danger of actually having to face his own ego and let go of it through real submission to spiritual discipline.

As far as I am concerned, all this is only an excuse to do nothing and take no spiritual path seriously while pronouncing ``infallible" judgements on whether this or that author is esoteric or not, or if Christianity, a religion for the masses, actually deserves even the slightest glance from such ``elitists" as themselves. ( And in this whole paragraph I have been referring not to the traditionalists as such, but to their pedantic followers who are plaguing the internet these days).

Like I said, I have profited and still profit immensely from the writings of Guenon ``and friends", but I do not slavishly follow them on every issue. When Guenon says one thing and a most reputed Father says something different, there can be no question regarding my allegiance- and this because I really take Guenon's work seriously.


\hfill

\texttt{Graham on 2013-06-02 at 23:55 said: }

Hopefully I will have more time to engage this essay soon. The first thing I find noteworthy is that Evola's evaluation of Catholicism here is more favourable, and far more nuanced, than that of some of his contemporary followers. Evola evidently deems Christianity rich in esoteric symbolism, connected at its origins with genuine initiation into the greater mysteries, and calls it a traditional form (if with qualification). This is all very far from any `strong Nietzschean' readings of his works.

Otherwise, I still think it's possible that Evola suffered a certain spiritual pride that affected his notion of deification, entwining it with a wicked fantasy of the `destruction of human nature' (and didn't Mircea Eliade remark the same thing?). Self-evidently, then, anything that smacked faintly of `sentimentalism' or `mere humanity’ – such as the theological virtues, the devotions of certain mystics and saints, or perhaps the Johannine symbolism of God as Love – was taken by him as proof of spiritual inferiority. While fully accepting the superiority of intellect to feeling, might this contempt not be the projection of a man, preoccupied with erasing his humanity, onto a type that instead sought its purification and integration – devotion, sentiment and all?


\hfill

\texttt{Matt on 2013-06-03 at 17:48 said: }

I've always considered Evola's statements on the destruction of human nature to be read in the Shaivite sense; human nature/human state isn't ``destroyed" but transcended.

Now deep down Evola might not have meant that, but at least from reading his writings, that seems to be what he meant.


\end{sffamily}\end{footnotesize}
